% =============================================================================
% §1.7 重写草稿 v8（独立片段，未 patch 进主稿；将替换主稿旧 §1.7 第 809–899 行）
% 相对 v7，依据 v7 独立专家复审报告 section_1_7_v7_review_report_20260617_zh.md
% （编辑部决定：Minor–Major Revision；0 CRITICAL，但查出 v6 轮系统复审漏掉的 3 处实质问题）
% 及与作者逐项定夺 + 作者独立二次核验，落实 5 个 MAJOR + 一批 MINOR：
% ---- MAJOR ----
%   [M1] 容量匹配 over-claim：实测 A1=M0=62,806（等参），A2=60,871（−3%），A4=62,230（−1%），
%        A3 no_lift=43,079（−31%，移除 J_θ 分支即减参）；黑箱才加宽（se3_accel 1.88/fullstate 1.78）。
%        → 把"各候选按参数量匹配"限定到跨族(PH↔黑箱)；消融诚实拆分（A1 干净等容量单开关；
%          A3 容量与作用不可分、结论条件化）。证据：各 best_model.pt state_dict 实测。
%   [M2] 海流"块内分段常值"误述数据生成：实测数据生成海流为连续缓变 OU（dt_sim=0.01 演化、
%        每 0.05s 快照 v_c_snap），块内常值的是控制命令 u_cmd；块内冻结(block_v_c)属"模型契约层"
%        （AUVHamNODE.py:119–120 "v_c^n carried with zero derivative over one integration window"）。
%        → 数据生成如实写连续 OU 外源；把块内冻结移到训练目标/rollout 的模型契约层。§1.6 不动。
%   [M3] 主指标只报种子均值、无离散度：summarize_sweep.py 已算种子间 mean/std/min/max（未算 IQR）。
%        → 主精度随附种子间 ±std；N=5 小样本、非 i.i.d.，只作描述性离散度、不做假设检验。
%   [M4] 缺积分步长 + 零敏感性：实测 odeint(method=rk4) 无 step_size → 步长=t_eval 间隔 0.05s；
%        数据生成 dt_sim=0.01s。仓库 200/200 config 全 rk4、无任何步长/求解器敏感性落档。
%        → 表中补步长 0.05s（数据生成 0.01s）；写设计理据（钉在观测间隔、三处共用）+ 诚实 limitation
%          （未做步长敏感性、结论限定 RK4@0.05s）。【红线：不得声称做过敏感性实验——无落档证据】。
%   [M5] 因果语言与自认混杂相互拉扯：把"分离/隔离各因素作用"降为"比较各结构假设诱导的长期行为差异"；
%        区分"显式容量(已跨族匹配)"与"结构归纳偏置(被测对象、非混杂)"；识别担子落等容量单开关 A1。
% ---- MINOR ----
%   m1 海流上限 0.5 是初值采样上限、OU 漂移后瞬时约 0.75（并入 M2）。
%   m2 M5 锚定标准 Neural ODE 并 cite chen2018neuralode（键存在于 auv_hamnode_refs.bib）。
%   m3 补"为何不用李群/流形积分"取舍句（几何保持作被检验量；并入 M4 数值实现段）。
%   m4 N 定义=对该指标产生有限运行级统计量的种子数（并入 M3）。
%   m5 60s 自由递推使状态推离训练分布(OOD)，有效性为被检验项。
%   m6 "长期预测可用性的一部分"→"长期预测稳健性的一个必要侧面(初值敏感性)"，并说明不做过程噪声的范围界定。
%   m7 "评估时域三层"→"一条 60s 递推在 10/30/60s 三时刻切片读取"（实测三 horizon completed_time 同为 59.80）。
%   m8 A4 同关两开关→合并为"输入条件化范围"单一较粗假设。
%   m9 机械能量诊断仅对含显式能量结构候选有定义（M4/M5 不适用）。
%   m10 激励质量评分=六自由度速度统计(std+极差)加权、横/角自由度更高权重（_excitation_score）。
%   m11 采样结构主实验默认逐控制块独立采样（整条共享=协议核查变体）。
%   m12 表"网络容量"补量纲锚点（隐藏宽度 128、M0 约 6×10⁴、黑箱按比例加宽）。
%   m13 "数值溢出"→"数值非有限(NaN/Inf)"（对齐 nan_or_inf；line 140/149/表 三处）。
%   m14 正文与两表数字/辅助指标逐字重复→正文留口径、清单让给表。
%   m15 §1.7 内重复声明(速度契约/生存条件量)二次出现改指回式收束。
%   m16 引号风格统一→留 patch 阶段（本片段不动）。
% ---- 红线（守住）----
%   §1.7 不报性能/排序/「最重要结构」（参数量属 setup 非性能、±std 是口径承诺非数值）；
%   三轴正交（训练协议{clean,iid,v4lite}/评估配置/评估协议）；
%   噪声=IC 鲁棒性正则非导航后验、预算为物理 DR 参照与数据集无关；
%   PH 仅开放六自由度机械核心、τθ≠G(q)u、海流仅在模型契约层分段常值、普通积分器不严格保持 SO(3)；
%   不宣称 v4lite 优、不预锁 §1.8 证据子集；M4 不得伪造步长敏感性实验。
%   label sec:training-baselines-protocol 保留。
% ---- patch 阶段连带（超出本片段）----
%   §1.8 同步 M0–M5/A1–A4 代号（每模型首现 + 四表加代号列）；
%   §1.8 离散度规划（本轮已锁形式）：四张结果表主精度列（60s 终端位置误差 中位数/P95）改 "值 ± std"
%     （种子间 std），完成率/失败原因行不加离散度；全发散模型(M5 型)仍记"长期稳定性失败"、不报中位数；
%     表注统一注明"± 为五种子间标准差，N 见各行"；§1.8 文件实际补列与引号统一留 patch 阶段。
% 真值源：data_collection.py（海流连续 OU、_excitation_score、dt_sim/dt_ctrl）、
%   AUVHamNODE.py:12–26,119–120（v_c^n 块内零导数契约）、train_utils.py:324,2230（rk4 无 step_size）、
%   rollout_benchmark_engine.py:222（block_v_c）、auv_model_registry.py、
%   scripts/summarize_sweep.py:95–115（种子间 mean/std/min/max）、各 best_model.pt state_dict（参数量实测）、
%   checkpoints/sweep_oc_phase1a_decision_clean_t2_wpfrag_phnode_full/.../config.json。
% =============================================================================

\section{实验设置}
\label{sec:training-baselines-protocol}

第~\ref{sec:structured-continuous-model} 节的参数化在函数类层面保证了正定逆质量、半正定耗散与斜对称零功率耦合，第~\ref{sec:energy-power-properties} 节给出了静水机械核心的功率平衡；这些性质由构造保证、不随训练改变。然而，这些结构约束是否真的转化为更可靠的长期运动预测，并不能由构造本身回答。模型作为长期轨迹预测器的实际行为，取决于数值积分、状态转换与数据分布的共同作用，只能通过训练后的长期递推与诊断来检验。本章实验因此围绕一个问题展开——在受控仿真数据与给定模型族下，逐项结构假设如何影响长期自由递推的稳定性与精度——并以一组在不同结构维度上受控变化的对照来回答。本节规定回答该问题所需的实验设置：数据如何生成（第~\ref{subsec:data-generation} 小节）、训练目标（第~\ref{subsec:training-objective} 小节）、用于比较各结构假设的对比模型与消融（第~\ref{subsec:models-ablations} 小节）、短期与长期评估方式（第~\ref{subsec:rollout-eval} 小节）、评价指标（第~\ref{subsec:metrics} 小节），以及作为鲁棒性检验的初值扰动（第~\ref{subsec:ic-perturbation} 小节）。各结构因素的定量作用在第~\ref{sec:results-structure-evidence} 节报告，本节不预先给出任何性能结论。

\subsection{数据生成}
\label{subsec:data-generation}

实验数据由 REMUS~100 自主水下航行器的六自由度刚体仿真器生成。仿真器以体坐标动力学为核心，采用显式时间积分推进平移与角运动，姿态以四元数表示以避免欧拉角奇异；推进器转速与舵面命令、海流与深度上下文作为外源输入驱动运动。每条轨迹由随机生成的命令序列激励，并切分为若干控制块：块内控制命令保持分段常值、块间更新，海流等外源流场则按仿真步连续演化，从而既覆盖丰富的受控机动，又与第~\ref{subsec:training-objective} 小节按控制块组织的训练目标一致。

命令序列的设计目标，是在有限轨迹内充分激励航行器动力学，使学到的连续时间向量场在尽可能宽的频段与工况下受到约束，而非只覆盖单一机动模式。为此，舵 \(\delta_r\)、艉舵 \(\delta_s\) 与推进器转速由三类互补的激励信号按固定比例混合生成：伪随机二进制序列（PRBS，占比 \(0.40\)）在随机保持时长后切换、经平滑后给出宽频的开关式激励；多正弦扫频（CHIRP，占比 \(0.35\)）以叠加的扫频正弦覆盖舵与艉舵的特征频段，刻画频率响应；Ornstein--Uhlenbeck 均值回复随机过程（OU，占比 \(0.25\)）给出平滑且时间相关的随机机动。推进器转速在 PRBS 与 OU 场景下由 OU 过程驱动、在 CHIRP 场景下取低频正弦，整体保持平滑变化。三类信号——开关式宽频、确定性扫频与平滑随机——在频率与幅度上互补，共同保证数据集对动力学的充分激励。海流作为连续缓变的外源流场注入：每条轨迹随机采样初始流速（在 \([0,0.5]~\mathrm{m/s}\) 内取值）与流向，叠加 Ornstein--Uhlenbeck 缓慢漂移并按仿真步连续演化，使航行器经历连续时变的流场（经漂移与截断后，瞬时流速上限约 \(0.75~\mathrm{m/s}\)）；有海流样本依式~\eqref{eq:velocity-contract} 的速度契约维持总体速度与相对水速度的转换。

生成的原始轨迹经两阶段质量过滤后入库：第一阶段剔除违反 \(\SO(3)\) 正交性或速度界限的物理无效轨迹；第二阶段按激励质量评分的百分位阈值保留信息量充足的样本（保留评分最高的约 \(70\%\)），使数据集偏向能有效约束动力学的受激轨迹，而非低激励的近平凡运动。这里的激励质量评分是各轨迹六自由度速度统计（速度标准差与极差）的加权度量，并对横向与角运动自由度赋予更高权重——鱼雷形 AUV 的这些自由度更难被充分激励（实现见 \texttt{data\_collection.py} 中的 \texttt{\_excitation\_score}）。需说明，该激励过滤使数据偏向高信息工况，其对各模型相对表现的潜在影响在本设置下未单独量化。有海流数据集即按上述流程生成、过滤并入库，每条轨迹切分为等长控制块；数据规模、控制块时长、状态采样间隔与训练--验证划分等定量设置统一列于表~\ref{tab:experimental-config}。仿真器可在有海流与无海流两种设置下生成数据；本章实验聚焦有海流环境，全部训练与评估均在该环境下进行。数据态在速度分量上保存体坐标系总体速度 \(\nu_b\)，训练前按式~\eqref{eq:velocity-contract} 转换为模型态相对水速度 \(\nu_r\)。数据生成与模型训练各自记录随机种子，跨模型比较固定同一数据镜像，使模型间差异不掺入数据采样差异。仿真器的水动力系数与执行器参数取自 REMUS~100 的公开数据，仅用于生成轨迹并初始化机械先验，不作为可从轨迹唯一辨识的真值（第~\ref{sec:modeling-prelim} 节）。

\subsection{训练目标}
\label{subsec:training-objective}

训练以控制块为基本单位（第~\ref{subsec:data-generation} 小节）：在单个控制块内，模型契约将控制命令与海流等外源上下文取为分段常值——即在该积分窗口内冻结于块初值、其导数为零——模型从块初值积分生成预测轨迹，
\begin{equation}
  y_0=\mathcal{T}_{d\to m}(s_0),\qquad
  \hat y(t_i)=\Phi_\theta(t_i;t_0,y_0),\qquad
  \hat s_i=\mathcal{T}_{m\to d}(\hat y(t_i)),
  \label{eq:block-prediction}
\end{equation}
其中 \(\mathcal{T}_{d\to m}\) 与 \(\mathcal{T}_{m\to d}\) 为式~\eqref{eq:data-to-model-state} 与式~\eqref{eq:model-to-data-state} 的数据态--模型态转换，\(\Phi_\theta\) 为式~\eqref{eq:complete-augmented-vector-field} 向量场经神经常微分方程积分得到的流映射。目标轨迹同样转换到模型态，使有海流样本的速度监督作用于相对水速度 \(\nu_r\)（依式~\eqref{eq:velocity-contract} 的速度契约），位置、姿态与执行器状态在两种表示中保持同一数值。需说明，数据真值的海流在控制块内仍按仿真步缓慢漂移（第~\ref{subsec:data-generation} 小节），而模型契约将其近似为该块内的分段常值；这一契约层简化使块级初值问题良定，海流的连续时变性则由块间更新与长期递推协议承接（第~\ref{subsec:rollout-eval} 小节）。

一个典型的加权训练目标可写为
\begin{equation}
  \mathcal{L}
  =
  \sum_i
  \lambda_x\norm{\hat x_i-x_i}^2
  +\lambda_R d_{\SO(3)}(\hat R_i,R_i)^2
  +\lambda_\nu\norm{\hat\nu_{r,i}-\nu_{r,i}}^2
  +\lambda_a\norm{\hat u_{a,i}-u_{a,i}}^2
  +\lambda_{\SO}\mathcal{L}_{\SO(3)} .
  \label{eq:training-loss}
\end{equation}
姿态误差以旋转测地距离 \(d_{\SO(3)}(\hat R,R)\) 度量、而非旋转矩阵的欧氏差，使监督直接作用于 \(\SO(3)\) 流形上的几何偏离，并对脱离旋转群的非物理偏移保持敏感；\(\mathcal{L}_{\SO(3)}\) 为旋转矩阵的正交性与行列式正则项。各分量残差在计入前按其经验标准差归一，使不同量纲的状态分量在目标中可比；默认位置、姿态与相对速度项等权（权重 \(1\)），执行器项权重取 \(0.2\)、\(\SO(3)\) 正则权重取 \(10^{-3}\)，连同 AdamW 权重衰减 \(10^{-4}\) 等具体取值记录于运行配置。速度项以相对水速度为监督对象，与第~\ref{sec:state-current-contract} 节的速度契约一致；评估时预测轨迹经 \(\mathcal{T}_{m\to d}\) 回到数据态，终端位置、姿态与总体速度等可观测量在数据空间解释（第~\ref{subsec:metrics} 小节）。连续时间积分、优化器与网络容量等实现细节统一列于表~\ref{tab:experimental-config}。

\begin{table}[htbp]
  \centering
  \caption{长期递推实验的主要配置}
  \label{tab:experimental-config}
  \begin{tabularx}{\textwidth}{>{\raggedright\arraybackslash}p{0.18\textwidth}>{\raggedright\arraybackslash}X}
    \toprule
    配置维度 & 设置 \\
    \midrule
    数据规模 & 海流数据集生成一千条原始轨迹、经两阶段过滤后入库，每条一百五十个控制块；控制块时长 \(0.2~\mathrm{s}\)，状态采样间隔 \(0.05~\mathrm{s}\)，过滤后样本按 \(0.8\) 划分训练与验证集（约 \(559/140\) 条）。 \\
    激励信号 & 舵、艉舵与转速由 PRBS、CHIRP、OU 三类激励按 \(0.40/0.35/0.25\) 混合生成；海流为连续缓变的 OU 外源流场。 \\
    速度契约 & 数据态保存体坐标系总体速度 \(\nu_b\)，训练前按式~\eqref{eq:velocity-contract} 转换为模型态相对水速度 \(\nu_r\)。 \\
    数值积分 & 固定步长四阶 Runge--Kutta，步长 \(0.05~\mathrm{s}\)（与状态采样间隔一致；数据生成仿真步为更细的 \(0.01~\mathrm{s}\)），单精度浮点状态；训练、块级评估与长期递推共用同一求解器与步长。 \\
    网络容量 & 跨族对照中黑箱基线主干网络宽度按可训练参数量加宽至与完整模型 M0 同量级（主干隐藏宽度 \(128\)，M0 约 \(6\times10^4\) 个可训练参数，黑箱按比例加宽）；结构消融按构造改变参数量（A1 与 M0 等参，A3 因移除零功率耦合分支约减少三成）。含质量结构的候选以 REMUS~100 质量矩阵与执行器先验初始化。 \\
    优化 & AdamW，学习率经线性预热后余弦衰减，至多三百个 epoch；批量、学习率端点与正则权重记录于运行配置。 \\
    种子 & 每个模型在五个固定随机种子（42--46）下独立训练与评估；跨种子聚合取运行级统计量的种子间均值，并随附种子间标准差。 \\
    递推评估 & 独立重采样初值（不与训练/验证集重叠），覆盖伪随机二进制序列（PRBS）、多正弦扫频（CHIRP）与 Ornstein--Uhlenbeck 过程（OU）三类控制场景，每场景 30 条评估轨迹；在同一条 60~\(\mathrm{s}\) 自由递推上，于 10、30 与 60~\(\mathrm{s}\) 三个时刻切片读取终端量，主时域 60~\(\mathrm{s}\)。 \\
    \bottomrule
  \end{tabularx}
\end{table}

\subsection{对比模型与结构消融}
\label{subsec:models-ablations}

结构约束的作用通过一组受控对照检验（表~\ref{tab:model-comparison-roles}）。这些候选按结构假设分组，而非排成单调的完整度阶梯：完整结构化模型 AUVHamNODE（M0）同时具备 \(\SE(3)\) 位姿几何、机械存储与 \(D_\theta\)/\(J_\theta\)/\(\tau_\theta\) 分解；端口哈密顿风格基线（M1、M2）保留机械存储核心，是两种平行的广义力组织方式——M1 以一般位形广义力替代标量势能诱导的保守力，M2 把非保守广义力合并而非分解；\(\SE(3)\) 动量黑箱（M3）保留位姿几何与常质量映射、仅将动量动力学交由黑箱，\(\SE(3)\) 加速度黑箱（M4）只保留位姿几何、把加速度动力学交由黑箱而不再保留显式质量映射；全状态黑箱（M5）则不含任何显式几何或能量结构，对应不含物理结构的标准神经常微分方程基线~\citep{chen2018neuralode}。各候选因而在不同结构维度上各移除或改组一类因素；对照所比较的，是各结构假设在相同数据与训练条件下诱导的长期行为差异，而非对真实水动力机制的因果识别。分组本身不预设结构越完整表现越好。其中 M1--M4 参照 \(\SE(3)\) 哈密顿及李群端口哈密顿神经常微分方程的思路~\citep{duong2021se3hamnode,duong2024liegroupphnode}，但均在统一的 AUV 状态契约、数据协议与评估方式下重新实现为受控候选，而非复刻文献中的原始模型。

为使跨族对照（结构化端口哈密顿候选与全状态黑箱之间）反映结构差异而非显式容量差异，黑箱基线的主干网络宽度按可训练参数量加宽至与完整模型 M0 同量级（具体取值见表~\ref{tab:experimental-config}）；含质量结构的候选统一以 REMUS~100 的质量矩阵与执行器先验初始化，黑箱基线无对应结构组件、故不注入该先验。这里须区分两类“容量”：其一是\emph{显式容量}，即可训练参数量，已在跨族对照中按上述方式匹配；其二是\emph{结构归纳偏置}，如机械存储、零功率耦合与位姿几何，它们正是本研究有意比较的被测对象，而非需要消除的混杂。就此而言，结构消融在显式容量上并不一致：A1（无质量先验）仅移除物理初始化、不改变网络架构，与 M0 参数量相等，是干净的等容量单开关；A2、A4 改变相应结构分支，但参数量与 M0 基本持平；A3（无零功率耦合）移除斜对称分支 \(J_\theta\) 本身即减少该分支的参数，可训练参数量较 M0 约低三成，故一旦 A3 表现劣化，其作用与该分支容量在本设置下不完全可分。因此，对各结构假设作用较干净的识别主要依赖显式容量基本可比的消融——尤以等参且不改变架构的 A1 最为可靠，涉及 A3 的结论则相应条件化。综合而言，结构先验本身会改变有效容量，这一混杂无法仅由参数量匹配完全消除，故下文结论一律以“当前模型族、数据协议与数值实现下”为条件。

\begin{table}[htbp]
  \centering
  \caption{对比模型与结构消融及其检验目标}
  \label{tab:model-comparison-roles}
  \begin{tabularx}{\textwidth}{l>{\raggedright\arraybackslash}p{0.30\textwidth}>{\raggedright\arraybackslash}X}
    \toprule
    代号 & 模型 & 检验目标 \\
    \midrule
    M0 & 完整结构化模型 AUVHamNODE & 含速度契约、\(\SE(3)\) 运动学、机械存储、\(D_\theta\)/\(J_\theta\)/\(\tau_\theta\) 分解与执行器滞后的完整候选项。 \\
    M1 & 端口哈密顿位形广义力基线 & 保留机械存储核心，以一般位形广义力替代标量势能诱导的保守力，检验能量核心的作用。 \\
    M2 & 端口哈密顿合并广义力基线 & 保留机械存储核心，把非保守广义力合并而非分解，检验广义力分解的作用。 \\
    M3 & \(\SE(3)\) 动量黑箱 & 保留位姿几何与常质量映射，动量动力学交由黑箱。 \\
    M4 & \(\SE(3)\) 加速度黑箱 & 保留位姿几何，加速度动力学交由黑箱、不保留显式质量映射。 \\
    M5 & 全状态黑箱 & 不含显式位姿几何与机械能量结构的非结构化状态导数模型，对应标准神经常微分方程基线。 \\
    \midrule
    A1 & 无质量先验消融 & 移除物理质量矩阵的先验初始化，保留正定质量参数化（与 M0 等参数）。 \\
    A2 & 对角阻尼消融 & 半正定耦合耗散矩阵 \(D_\theta(\xi)\) 限制为非负对角形式。 \\
    A3 & 无零功率（升力）耦合消融 & 移除斜对称零功率分支 \(J_\theta(\xi)\nu_r\)（同时减少该分支参数，约 \(-31\%\)）。 \\
    A4 & 广义力条件化消融 & 广义力分支 \(\tau_\theta\) 收缩为主要由执行器状态或命令驱动，合并速度条件与海流特征为单一输入条件化范围。 \\
    \bottomrule
  \end{tabularx}
\end{table}

结构消融（A1--A4）在完整模型 M0 上逐项移除或弱化单个结构组件（表~\ref{tab:model-comparison-roles}），每个消融对应一条可检验的结构假设，并与第~\ref{sec:energy-power-properties} 节的相应能量--功率性质相对接：A1 检验物理质量先验对动量--速度配对学习与长期递推的作用；A2 检验跨通道耗散耦合相对独立通道耗散的作用；A3 检验零功率升力类耦合在候选模型族内的作用（如前所述，移除该分支同时减少其参数，故其作用与分支容量不完全可分，亦不等同于断言真实水动力升力项被唯一辨识）；A4 将开放广义力通道的条件化范围收缩为主要由执行器状态或命令驱动，从而把“速度条件”与“海流特征”两类输入合并为单一较粗的条件化假设而非分别考察，检验广义力通道输入条件化范围的作用。

所有候选在统一的数据、训练与数值实现配置下比较（表~\ref{tab:experimental-config}），命令级参数与运行元数据保留在各运行记录中。进入定量比较的运行须共享同一数据协议、代码环境与评估方式；出现训练崩溃或长期递推发散的运行移出精度聚合并单独披露，其中全部种子均发散者不报告中位数、而记为长期稳定性失败，以与精度不佳相区别。异常的判定与披露规则见第~\ref{sec:results-structure-evidence} 节。

\subsection{评估方式：短期与长期}
\label{subsec:rollout-eval}

块级训练目标只反映局部受控初值问题的拟合质量，无法说明误差在自由递推中的累积；因此评估分短期与长期两个层次，二者考核不同的能力。

短期评估在留出的控制块与轨迹上度量模型对局部动力学的逼近：以留出样本的块初值启动，积分单个控制块，比较预测与真值的状态偏差。短期评估反映模型是否学到了正确的瞬时状态演化，是长期能力的必要而非充分条件——块级精度高的模型仍可能在自由递推中发散。

长期评估为本章主要评估方式，采用控制块串接的自由递推：首块由给定初值启动，其后每块以上一块的预测末状态为初值，并按递推契约逐块更新位置原点、海流速度、控制命令与深度上下文；与训练时一致，海流在每个控制块内取常值、块间更新。所得轨迹不再由真实状态校正，单块误差沿时间累积、放大，姿态漂移、速度发散与数值失败等仅在长期显现的问题由此暴露。长期评估使用受控仿真生成的有海流 AUV 轨迹；有海流数据依式~\eqref{eq:velocity-contract} 维持总体速度与相对水速度的转换，使位姿几何按总体速度推进、水动力配对作用于相对水速度。评估轨迹由独立重采样生成，与训练、验证集不重叠，以避免数据泄漏。

评估时域并非三次独立的递推，而是在同一条自由递推上读取：以 60~\(\mathrm{s}\) 为本章主时域，并在 10、30 与 60~\(\mathrm{s}\) 三个时刻切片读取终端量；较短时刻用于刻画误差随时间的增长形态与提前失败的位置，避免单一终端时刻掩盖中途已发散的轨迹。需说明，单条训练轨迹长 30~\(\mathrm{s}\)（一百五十个控制块、每块 \(0.2~\mathrm{s}\)），而长期递推评估串接至 60~\(\mathrm{s}\)：训练目标定义在单个控制块上、与轨迹总长解耦，递推协议可将任意数目的控制块串接，故 60~\(\mathrm{s}\) 主时域考核的是超出单条训练轨迹长度的累积行为；模型在该长度上的稳定性是被检验项，而非由训练目标保证的性质。还应指出，自由递推使预测状态逐步偏离训练数据所覆盖的分布，60~\(\mathrm{s}\) 段的有效性因而取决于递推状态是否仍落在训练分布的支撑之内——这一点同样是被检验的，而非被假定的。未完成目标时域的轨迹，其终端误差不计入该时域的条件误差统计；因此该条件终端误差是在各模型各自完成的轨迹子集上度量的生存条件量，模型间完成率显著不同时这些子集并不一致，其横向比较须结合完成率解读（第~\ref{subsec:metrics} 小节）。对未完成的轨迹，进一步区分模型预测发散、真值侧发散、数值求解失败与数值非有限（NaN/Inf），使后续分析能把结构性不稳定与偶发的数值失败分开归因。

\subsection{评价指标}
\label{subsec:metrics}

评价指标分两类，分别回答“模型预测得准不准”和“结构组件是否按设计工作”。任务精度指标在数据空间定义，刻画模型作为轨迹预测器的可观测性能，模型优劣以其判定；物理诊断量在模型空间定义，核验各结构组件的行为是否与设计一致，但不替代任务精度。短期与长期评估各取其中适用者：短期以留出块上的逐步 RMSE 度量局部拟合，长期以自由递推的终端与轨迹聚合误差度量误差累积。完整指标集列于表~\ref{tab:metrics}，下文说明主指标的取舍。

长期递推的主精度取 60~\(\mathrm{s}\) 终端位置误差的中位数与第 95 百分位（P95）并列度量。之所以以终端位置误差为主，是因为它是自由递推中误差累积的最终可观测后果，并在数据空间有直接的任务意义（航位精度）；之所以中位数与 P95 并列而非单一标量，是因为典型精度与尾部稳健并不总是同向——一个中位数很低的模型仍可能在少数轨迹上累积显著更大的误差，单一中位数或均值都会掩盖这一点。若出现典型精度与尾部排序不一致的情形，并列双指标即可分别反映，相关情形见第~\ref{sec:results-structure-evidence} 节。除主指标外，本章还报告终端深度误差（位置误差的垂向分量，单列以反映 AUV 定深任务的工程意义）、旋转测地误差与线、角速度误差等分通道辅助精度指标，刻画位置以外各状态分量的预测质量，其中线速度同时给出相对水速度与总体速度两种表示，以核验速度契约（第~\ref{sec:state-current-contract} 节）；短期评估则以位置 RMSE、旋转测地误差与速度 RMSE（含 \(u,v,w\) 与 \(p,q,r\) 分量）度量块内逐步拟合，回答模型是否学到正确的瞬时演化。

物理诊断量核验结构组件是否按设计工作：相对水速度误差核验 \(\nu_r\) 是否与速度契约一致；机械能量跨度与最大能量变化反映机械存储是否在递推中保持有限量级，这一诊断仅对含显式能量结构的候选有定义（不含能量结构的 M4、M5 等不适用）；\(\SO(3)\) 行列式与正交性误差反映姿态矩阵是否保持接近旋转群。

这里须说明本章对数值积分的处理及其取舍。所有模型均采用通用的固定步长积分器，步长取 \(0.05~\mathrm{s}\)、与状态观测间隔一致（见表~\ref{tab:experimental-config}），使训练、块级评估与长期递推在同一离散化下进行，不引入步长口径上的差异。这一选择是有意的：通用积分器并不严格保持 \(\SO(3)\)，因而几何保持（如正交性偏离）在本章被当作可被诊断、并可由结构隐式改善的\emph{被检验量}，而非由积分器从外部强制施加为零的约束；这与本章“以受控对照检验结构作用”的取向一致。同时须如实指出，本章未单独考察积分步长或求解器阶数的敏感性，因此涉及长期稳定性的结论均限定于该固定步长四阶 Runge--Kutta、步长 \(0.05~\mathrm{s}\) 的数值设置之下。

完成率与失败原因不参与精度排序，但决定终端误差是否具可比性：如前所述，终端误差只在完成目标时域的样本上度量，故为生存条件量；当被比较模型的完成率或有效种子数 \(N\) 存在显著差异时，终端中位数不可单独横比，须与完成率、\(N\) 联合解读，全部种子发散者不报告中位数而记为长期稳定性失败。失败原因区分模型预测发散、真值侧发散、数值求解失败与数值非有限（NaN/Inf）；深度越界则作为约束违反率单独统计，与上述失败原因分类并行记录（二者可同时出现），使后续分析能把结构性不稳定与偶发的数值失败分开归因。

各精度量先在单个运行内对评估轨迹取统计量，再在随机种子之间取均值，并报告种子间标准差以反映跨种子离散度；记号 \(N\) 指对该指标产生有限运行级统计量的有效种子数。鉴于 \(N=5\) 属小样本、且各种子结果并非独立同分布，本章只以种子间标准差作描述性的离散度报告，不进行正式的假设检验。

\begin{table}[htbp]
  \centering
  \caption{短期与长期评估的指标体系}
  \label{tab:metrics}
  \begin{tabularx}{\textwidth}{>{\raggedright\arraybackslash}p{0.13\textwidth}>{\raggedright\arraybackslash}X>{\raggedright\arraybackslash}X}
    \toprule
    类别 & 短期评估（留出块/轨迹） & 长期评估（自由递推） \\
    \midrule
    任务精度 & 位置 RMSE、旋转测地误差、线/角速度 RMSE 及 \(u,v,w\)/\(p,q,r\) 分量误差 & 终端位置误差中位数与 P95（主）；终端深度、旋转测地与相对/总体线、角速度误差；轨迹平均/最大位置与姿态误差；块级局部误差 \\
    物理诊断 & \(\SO(3)\) 行列式与正交性误差 & \(\SO(3)\) 正交性误差（最大值）、机械能量跨度、最大能量变化、终端机械能量 \\
    完成性 & 块成功率、求解失败计数 & 完成率、失败原因（模型预测发散/真值发散/求解器失败/数值非有限）、深度越界违反率 \\
    \bottomrule
  \end{tabularx}
\end{table}

\subsection{初值扰动与鲁棒性评估}
\label{subsec:ic-perturbation}

前述评估默认初值精确已知。但在部署中，模型实际使用的初始状态来自导航系统对航位推算与惯性测量的滤波估计，必然带有误差；评估模型在带噪初值下的稳健性，因而是长期预测稳健性的一个必要侧面，即对初值误差的敏感性。初值扰动正为此设计：它只对 \(t=0\) 的初始状态施加误差，模拟滤波后状态估计的残差，其后递推过程不再注入噪声。这一界定也划定了本章的考察范围：本章聚焦初值（滤波后状态）误差这一侧面，而过程噪声与全程观测噪声属于另一研究轴、不在本章范围；如此可使带噪初值下的差异较纯粹地归因于模型对初值误差的敏感性，而非过程噪声的累积。需强调，该接口刻画的是模型实际使用变量的误差半径，而非完整导航误差后验——它不复现导航滤波器的联合协方差结构，只约束初值各通道的误差预算。

初值扰动在训练与评估两个阶段引入，承担两个不同角色，分析中不应混为一谈：训练阶段，带噪初值作为鲁棒性正则，决定模型是否在带噪初值上学习；评估阶段，带噪初值作为鲁棒性测试，决定模型在何种扰动下受考核。同一种扰动既可作为训练干预、也可作为评估干预，其含义由所处阶段决定。

扰动的设计分三个维度。其一是受扰通道：噪声由干净数据态转为干净模型态后，在相对水速度 \(\nu_r\)、姿态 \(R\) 与执行器反馈 \(u_a\) 通道上采样零均值扰动，再启动递推；如此控制的是模型实际使用的变量，并在海流设置下保持 \(R\)、\(\nu_r\) 与海流的语义一致，避免数据态扰动看似不大、而模型态输入已被过度污染的失配。海流通道 \(v_c^n\) 是否扰动取决于所采用的导航参考（见下）。其二是扰动强度：评估按强度递增分为名义与退化两档，另设一档在名义随机扰动上叠加固定航向偏置的有偏配置，用以考核对系统性航向误差的稳健性。其三是采样结构：噪声可逐控制块独立采样，也可由同一轨迹的各控制块共享一次实现，二者在训练与评估两端均可实例化；本章带噪设置的主实验默认采用逐控制块独立采样，整条轨迹共享一次实现的变体用于协议层面的稳健性核查。

各通道预算以 REMUS~100 的航位推算（dead reckoning）传感器量级为物理参照、与训练数据分布无关。相对水速度噪声按“固定底噪加速度比例项”建模，
\begin{equation}
  \sigma_{\nu,i}=\sqrt{\,(\sigma_i^{\mathrm{floor}})^2+\big(k_i\,\lvert\nu_{r,i}\rvert\big)^2\,},
  \label{eq:nu-noise-model}
\end{equation}
其中 \(\sigma_i^{\mathrm{floor}}\) 为传感器底噪下限、\(k_i\) 为相对瞬时速度的比例系数——这一形式对应多普勒计程仪“按测速百分比加厘米每秒底噪”的精度规格；角速度与姿态噪声取固定预算。本章主实验采用标准航位推算参考（多普勒计程仪、罗经与速率陀螺），其预算列于表~\ref{tab:noise-budget}；在该参考下，海流不作为被独立估计的初值状态、其初值通道不施扰动。另保留一套高精度惯导（high-grade INS）参考作为扩展，其姿态与航向误差更小、并对海流通道引入非零预算。由于预算以物理传感器量级而非数据集统计定义，同一档位在不同数据集上的物理含义保持一致；各运行仍记录最终采用的预算以备追溯。

\begin{table}[htbp]
  \centering
  \caption{初值扰动各通道的噪声预算（REMUS~100 航位推算参照，本章主实验）}
  \label{tab:noise-budget}
  \small
  \begin{tabularx}{\textwidth}{>{\raggedright\arraybackslash}X >{\centering\arraybackslash}p{0.105\textwidth} >{\centering\arraybackslash}p{0.105\textwidth} >{\centering\arraybackslash}p{0.105\textwidth} >{\centering\arraybackslash}p{0.165\textwidth}}
    \toprule
    通道（参数） & 名义训练 & 名义评估 & 退化评估 & 航向偏置评估 \\
    \midrule
    相对线速度 底噪 \(\sigma^{\mathrm{floor}}\)（m/s，水平/垂向）& 0.002 / 0.003 & 0.003 / 0.004 & 0.006 / 0.008 & 同名义评估 \\
    相对线速度 速度比例 \(k\) & 0.002 & 0.003 & 0.006 & 同名义评估 \\
    相对角速度 std（rad/s，横滚俯仰/偏航）& 0.0006 / 0.0010 & 0.0008 / 0.0015 & 0.0015 / 0.0030 & 同名义评估 \\
    姿态 std（rad，横滚俯仰/偏航）& 0.0015 / 0.0060 & 0.0020 / 0.0100 & 0.0040 / 0.0250 & 名义评估 \(+\) 偏航偏置 0.035 \\
    海流 \(v_c^n\)（m/s）& 0（不施扰动）& 0 & 0 & 0 \\
    执行器 \(\delta_r,\delta_s\)（rad）/ 转速（r/min）& 0.002 / 3 & 0.003 / 5 & 0.004 / 10 & 同名义评估 \\
    \bottomrule
  \end{tabularx}
  \par\vspace{2pt}
  \footnotesize 注：相对线速度按式~\eqref{eq:nu-noise-model} 取值，底噪与比例系数源自多普勒计程仪精度规格；角速度与姿态取固定预算。海流通道在本章的航位推算参考下不施扰动（高精度惯导扩展参考下引入非零海流预算）。航向偏置评估在名义评估的随机扰动上叠加固定 \(0.035~\mathrm{rad}\) 的偏航偏置，符号按样本固定、一条递推内不变。采样结构（逐控制块独立采样或整条轨迹共享一次实现）由运行配置控制。
\end{table}

训练端对噪声强度采用预热--爬坡--混合的渐进调度，以避免噪声初值在优化初期破坏训练稳定性，具体预热轮数、爬坡区间与混合比例记录于运行配置；评估端对所有模型固定噪声随机种子，使被扰动初值在不同模型之间可直接比较。
